\chapter{Doped-Silicon Parent Systems and Compatibility}
\index{silicon!doped program|textbf}

\label{ch:doped-silicon-program}

\chapterstatus{expository synthesis of the accepted phosphorus and boron
branches; no production dopant calculation or impurity operator is claimed.}

The dopant program begins only after the bulk representation and its alignment
conventions have been identified. It treats phosphorus and boron as distinct
physical branches rather than interchangeable instances of an abstract defect.

\section{Compatible host and dopant branches}
\index{compatibility!host and dopant}
\index{supercell!matched pristine}


Each doped calculation uses a periodic supercell built from the frozen host
lattice vectors, with one substitutional dopant and a matched pristine
comparison supercell. In the primary dopant branches, internal atomic
coordinates are relaxed while the approved bulk lattice vectors remain fixed;
unrelaxed structures are retained only as diagnostics. Geometry and site
correspondence must be identified before pristine and doped operators are
compared.

The immediate non-spin-polarized, non-SOC bulk pilot does not by itself supply
every host reference required by the dopant program. A phosphorus comparison
requires a compatible bulk reference with documented spin and SOC conventions.
Final boron acceptor and continuum claims require a matching fully relativistic,
noncollinear SOC bulk reference. Operators from these different spin spaces are
not directly subtractable without an explicit identification.

\section{Substitutional phosphorus}
\index{phosphorus!substitutional}
\index{donor}


The phosphorus branch replaces one silicon atom at a diamond-lattice site by
phosphorus in a periodic supercell. Its primary version-one charge state is
neutral $\mathrm{P}_{\mathrm{Si}}^{0}$. With the accepted pseudopotential
valences, this system has one additional valence electron relative to the
matched pristine silicon supercell and therefore an odd valence-electron count.
The primary branch is consequently scalar relativistic, collinearly spin
polarized, and non-SOC.

The neutral supercell contains both the donor-core perturbation relative to
silicon and the additional electron that occupies and screens that
perturbation. It must not be relabelled silently as an isolated ionized donor.
A charged $\mathrm{P}_{\mathrm{Si}}^{+}$ calculation, or an extraction intended
to isolate an ionized-core potential, requires a separately specified
controlled branch.

The eventual phosphorus questions concern donor energies relative to a
compatible conduction-band reference, valley composition, localization or
envelope diagnostics, and the spatial structure of an aligned impurity
operator. Each finite supercell represents a periodic dopant array. An
isolated-donor interpretation requires supercell-size evidence rather than the
presence of a single substituted atom in the simulation cell.

\citationtodo{medium priority}{Check the periodic-dopant-array and dilute-limit
warning against defect-supercell literature. Candidate: Freysoldt et al.
(2014), bibliography key \texttt{freysoldt2014} (already present).}

\section{Substitutional boron}
\index{boron!substitutional}
\index{acceptor}
\index{spin--orbit coupling}


The boron branch replaces one silicon atom by boron and uses neutral
$\mathrm{B}_{\mathrm{Si}}^{0}$ as its primary version-one charge state. The
accepted valence convention gives one fewer valence electron than in the
matched pristine supercell and again produces an odd electron count.

For final acceptor and continuum claims, the physical specification requires a
fully relativistic, noncollinear spinor calculation with spin--orbit coupling.
A scalar-relativistic non-SOC calculation may be retained only as an explicitly
labelled method-development branch. This distinction prevents an early
simplification from being presented as the final physical treatment of the
spin--orbit-coupled valence manifold.

\citationtodo{medium priority}{Check the requirement to preserve the
spin--orbit-coupled valence manifold in an acceptor continuum treatment against
the shallow-acceptor literature. Candidate: Baldereschi--Lipari (1973),
provisional key \texttt{baldereschiLipari1973}; the candidate record is present
for this editorial note.}

The eventual boron questions concern acceptor energies relative to a
compatible valence-band reference, valence-subspace fidelity, localization or
envelope diagnostics, and the radial structure of the aligned impurity
operator. Charged $\mathrm{B}_{\mathrm{Si}}^{-}$ calculations are deferred to
a separate controlled branch.

\section{Dopant research questions}
\index{dopant!research questions}


The dopant program asks which identifications are required before pristine and
doped operators can be compared, which nested lattice classes preserve the
aligned perturbation, and whether a continuum effective-mass description is
adequate for a declared donor or acceptor use. Answers require retained
calculations and separately declared acceptance rules.

\section{Interpretive limits}

A finite periodic supercell is not an isolated impurity. A Kohn--Sham band gap
is not an exact quasiparticle gap. Spectral agreement is not operator equality.
Matrix subtraction is not physically meaningful before common-space and
energy alignment. Software execution and passing tests establish neither
numerical convergence nor scientific validation unless the corresponding
evidence and acceptance rules are present.

These limits are part of the scientific problem rather than editorial caveats:
they determine which comparisons are well posed and which conclusions the
resulting hierarchy can support.

\citationtodo{high priority}{Check the Kohn--Sham versus quasiparticle-gap
warning against Perdew--Levy (1983) and Sham--Schluter (1983), bibliography keys
\texttt{perdewLevy1983} and \texttt{shamSchluter1983} (already present).}
